\section{\texorpdfstring{$\infty$-toposy}{Toposy wyższe} i snopy wyższe}\label{cat:sec:infty-topoi}

W rozdziale~\ref{cat:sec:grothendieck-topologies} rozpatrywaliśmy presnopy i snopy zbiorów, w rozdziale~\ref{cat:sec:toposes} ich kategorie, a w rozdziale~\ref{cat:sec:descent-stacks} zstępowanie obiektów wraz z izomorfizmami. Przejdziemy teraz do presnopów przestrzeni, czyli typów homotopijnych, w których oprócz punktów występują drogi, homotopie dróg i dalsze dane wyższych rzędów. W miejsce zwykłego warunku sklejania użyjemy równoważności przestrzeni odwzorowań. Zasadniczym narzędziem konstrukcji będzie teoria dostępnych lokalizacji refleksyjnych z rozdziału~\ref{cat:sec:presentable-infty-bousfield}. Nie będziemy zakładać, że dowolny $\infty$-topos jest kategorią snopów dla zadanej klasycznej topologii Grothendiecka: definicja za pomocą lokalizacji obejmuje szerszy problem.

\subsection{Presnopy przestrzeni i wyższe twierdzenie Yonedy}

Przez $\mathcal S$ oznaczamy $\infty$-kategorię przestrzeni, modelowaną przez kompleksy Kana z ich przestrzeniami odwzorowań. Przestrzeń dyskretna odpowiada zbiorowi, natomiast przestrzeń jednopunktowa $*$ jest końcowa, a przestrzeń pusta $\varnothing$ początkowa. Ustalamy uniwersa tak, aby kategoria indeksująca była mała, a rozpatrywana kategoria presnopów mogła być duża.

\begin{definicja}[Presnop przestrzeni]\label{cat:def:infty-presheaf}
Dla małej $\infty$-kategorii $\mathcal C$ kładziemy
\[
 \mathcal P(\mathcal C)=\operatorname{Fun}(\mathcal C^{\mathrm{op}},\mathcal S).
\]
Obiekty tej kategorii nazywamy presnopami przestrzeni. Dla $c\in\mathcal C$ przez $y(c)$ oznaczamy presnop reprezentowalny
\[
 y(c)(d)=\operatorname{Map}_{\mathcal C}(d,c).
\]
Włożenie Yonedy $y:\mathcal C\to\mathcal P(\mathcal C)$ wykorzystuje funktorialność przestrzeni odwzorowań w obu argumentach.
\end{definicja}

Jeżeli $\mathcal C=N\mathcal A$ dla zwykłej małej kategorii $\mathcal A$, to $y(c)$ przyjmuje wartości w przestrzeniach dyskretnych odpowiadających zbiorom $\operatorname{Hom}_{\mathcal A}(d,c)$ (przykład~\ref{cat:ex:mapping-ordinary}). Ogólny presnop na $N\mathcal A$ nie musi mieć dyskretnych wartości. Nie wystarczy także zadanie jego wartości na obiektach i klas homotopii odwzorowań restrykcji: musimy zachować wszystkie dane funktora do $\mathcal S$.

\begin{twierdzenie}[Wyższe twierdzenie Yonedy; wynik zewnętrzny]\label{cat:thm:infty-topos-yoneda}
Dla $c\in\mathcal C$ oraz $F\in\mathcal P(\mathcal C)$ istnieje naturalna równoważność
\begin{equation}\label{cat:eq:infty-topos-yoneda}
 \operatorname{Map}_{\mathcal P(\mathcal C)}(y(c),F)\simeq F(c).
\end{equation}
W szczególności $y$ jest pełnym i wiernym funktorem w sensie przestrzeni odwzorowań.
\end{twierdzenie}

Korzystamy z wyższego twierdzenia Yonedy z J.~Lurie, \emph{Higher Topos Theory}, \S5.1.3. Porównanie odpowiada ewaluacji transformacji naturalnej w identyczności obiektu $c$: z $\eta:y(c)\to F$ otrzymujemy punkt $\eta_c(\id_c)$ przestrzeni $F(c)$. W przeciwnym kierunku punkt $x\in F(c)$ wyznacza rodzinę odwzorowań $\operatorname{Map}_{\mathcal C}(d,c)\to F(d)$ przez działanie restrykcji na $x$. To wyjaśnia konstrukcję porównania; pełna równoważność obejmuje także rodziny parametryzowane przestrzeniami i wszystkie ich koherencje, co stanowi treść wyniku zewnętrznego. Podstawiając $F=y(e)$, otrzymujemy
\[
 \operatorname{Map}_{\mathcal P(\mathcal C)}(y(c),y(e))
 \simeq \operatorname{Map}_{\mathcal C}(c,e).
\]
Równoważność ta jest odwrotna do odwzorowania indukowanego przez $y$, gdyż ewaluacja transformacji odpowiadającej morfizmowi $c\to e$ na $\id_c$ odzyskuje ten morfizm. To daje pełną wierność.

\begin{twierdzenie}[Granice, kogranice i prezentowalność presnopów; wynik zewnętrzny]\label{cat:thm:infty-presheaf-presentable}
Kategoria $\mathcal P(\mathcal C)$ jest prezentowalna. Granice i kogranice małych diagramów oblicza się punktowo:
\[
 (\varprojlim_i F_i)(c)\simeq\varprojlim_i F_i(c),
 \qquad
 (\varinjlim_i F_i)(c)\simeq\varinjlim_i F_i(c).
\]
Każdy morfizm presnopów jest równoważnością wtedy i tylko wtedy, gdy jest równoważnością po ewaluacji w każdym obiekcie $c$.
\end{twierdzenie}

Źródło: \emph{Higher Topos Theory}, \S5.1.2 i \S5.5.3 (kategorie funktorów do kategorii prezentowalnych). Przestrzenie mają wszystkie małe granice i kogranice, więc dla ustalonego $c$ istnieją obiekty po prawej stronie obu wzorów. Twierdzenie o punktowym obliczaniu granic i kogranic w kategoriach funktorów zapewnia ich funktorialność oraz własność uniwersalną jako presnopów, a nie tylko istnienie osobnych granic przestrzeni. Kryterium równoważności jest zastosowaniem twierdzenia~\ref{cat:thm:quasicat-pointwise-equivalence} do kategorii funktorów. Prezentowalność wymaga małości $\mathcal C$ i prezentowalności $\mathcal S$; korzystamy z wyniku zewnętrznego, nie wyprowadzamy jej z samego istnienia kogranic.

\begin{stwierdzenie}[Reprezentowalne presnopy jako rodzina generatorów]\label{cat:prop:infty-presheaf-generators}
Rodzina reprezentowalnych presnopów wykrywa równoważności w $\mathcal P(\mathcal C)$. Każdy reprezentowalny presnop jest zwarty, tzn. $\aleph_0$-zwarty w sensie definicji~\ref{cat:def:infty-kappa-compact}.
\end{stwierdzenie}
\begin{proof}
Dla morfizmu $f:F\to G$ naturalność~\eqref{cat:eq:infty-topos-yoneda} identyfikuje odwzorowanie
\[
 \operatorname{Map}(y(c),F)\longrightarrow\operatorname{Map}(y(c),G)
\]
z $f(c):F(c)\to G(c)$. Jeśli wszystkie te odwzorowania są równoważnościami, to $f$ jest równoważnością przez twierdzenie~\ref{cat:thm:infty-presheaf-presentable}. Odwrotna implikacja wynika z funktorialności przestrzeni odwzorowań.

Dla dowolnego małego diagramu $(F_i)$, a więc także dla diagramu filtrowanego, mamy
\[
\begin{aligned}
 \operatorname{Map}(y(c),\varinjlim_i F_i)
 &\simeq(\varinjlim_i F_i)(c)\\
 &\simeq\varinjlim_i F_i(c)\\
 &\simeq\varinjlim_i\operatorname{Map}(y(c),F_i).
\end{aligned}
\]
Pierwsza i trzecia równoważność wynikają z twierdzenia Yonedy, druga z punktowego obliczania kogranic. Z naturalności tych równoważności wynika, że ich złożenie w kierunku odwrotnym jest właśnie porównaniem z definicji zwartości. Otrzymujemy zatem zwartość, a nawet zachowywanie wszystkich małych kogranic przez funktor $\operatorname{Map}(y(c),-)$.
\end{proof}

\begin{przyklad}[Presnopy na kategorii jednopunktowej i kategorii strzałki]\label{cat:ex:infty-presheaf-simple}
Dla kategorii końcowej $*$ mamy $\mathcal P(*)\simeq\mathcal S$: ewaluacja w jedynym obiekcie jest równoważnością z odwrotnością przypisującą przestrzeni stały presnop. Dla $\mathcal C=N[1]$, gdzie $[1]$ ma jeden nietrywialny morfizm $0\to1$, presnop jest diagramem $F(1)\to F(0)$ przestrzeni. Przykładowo $y(0)$ ma wartość $*$ w $0$ i $\varnothing$ w $1$, natomiast $y(1)$ ma wartość $*$ w obu obiektach. Wynika to z wyliczenia zbiorów morfizmów do $0$ i do $1$. Naturalność Yonedy daje $\operatorname{Map}(y(0),F)\simeq F(0)$ oraz $\operatorname{Map}(y(1),F)\simeq F(1)$; w drugim przypadku wartość w $0$ jest wyznaczona przez restrykcję wybranego punktu z $F(1)$, wraz z koherentnymi danymi transformacji.
\end{przyklad}

\subsection{Definicja toposu wyższego}

\begin{definicja}[$\infty$-topos]\label{cat:def:infty-topos}
$\infty$-kategorię $\mathcal X$ nazywamy $\infty$-toposem, jeśli istnieje mała $\infty$-kategoria $\mathcal C$ i dostępna lokalizacja refleksyjna
\[
 L:\mathcal P(\mathcal C)\rightleftarrows\mathcal X:i
\]
taka, że $L$ zachowuje skończone granice. Oznacza to, że $L\dashv i$, prawy sprzężony $i$ jest pełny i wierny oraz dostępny, a lokalizacja jest lewostronnie dokładna w sensie definicji~\ref{cat:def:infty-bousfield-left-exact}.
\end{definicja}

Jest to definicja~6.1.0.4 z \emph{Higher Topos Theory} (J.~Lurie, Princeton University Press, 2009, \url{https://www.math.ias.edu/~lurie/papers/HTT.pdf}). W tym ujęciu topologia Grothendiecka nie stanowi części danych definicji. Wymagamy, aby $\mathcal X$ była równoważna takiej kategorii lokalnej, ale nie wyróżniamy jednej prezentacji. Skończone granice obejmują obiekt końcowy i homotopijne iloczyny włókniste, a zatem dokładność z lewej jest warunkiem silniejszym od samego istnienia refleksji.

\begin{stwierdzenie}[Podstawowe własności prezentacji]\label{cat:prop:infty-topos-presentation}
Niech $L:\mathcal P(\mathcal C)\rightleftarrows\mathcal X:i$ będzie prezentacją $\infty$-toposu. Wówczas $\mathcal X$ jest prezentowalna, jej małe granice obliczają się po włożeniu przez $i$, a dla dowolnego małego diagramu $(X_j)$ w $\mathcal X$ mamy
\[
 \varinjlim\nolimits_{\mathcal X} X_j
 \simeq L(\varinjlim\nolimits_{\mathcal P(\mathcal C)}iX_j).
\]
Ponadto rodzina obiektów $L(y(c))$ wykrywa równoważności w $\mathcal X$.
\end{stwierdzenie}
\begin{proof}
Przez twierdzenie~\ref{cat:thm:infty-presheaf-presentable} kategoria $\mathcal P(\mathcal C)$ jest prezentowalna. Dostępność lokalizacji i stwierdzenie~\ref{cat:prop:infty-bousfield-small-generation} dają mały zbiór morfizmów $S$, względem którego obraz istotny $i$ jest kategorią obiektów lokalnych. Twierdzenie~\ref{cat:thm:infty-bousfield-existence} zapewnia prezentowalność tej kategorii, a więc $\mathcal X$. Obiekty lokalne są domknięte na granice w kategorii otaczającej przez stwierdzenie~\ref{cat:prop:infty-bousfield-closure}, toteż pełna wierność $i$ pozwala obliczać granice w kategorii presnopów. Dla kogranicy i dowolnego $T\in\mathcal X$ otrzymujemy
\[
\begin{aligned}
 \operatorname{Map}_{\mathcal X}
 (L(\varinjlim_j iX_j),T)
 &\simeq\operatorname{Map}_{\mathcal P(\mathcal C)}
  (\varinjlim_j iX_j,iT)\\
 &\simeq\varprojlim_j\operatorname{Map}_{\mathcal P(\mathcal C)}(iX_j,iT)\\
 &\simeq\varprojlim_j\operatorname{Map}_{\mathcal X}(X_j,T).
\end{aligned}
\]
Pierwsza równoważność pochodzi ze sprzężenia, druga z własności uniwersalnej kogranicy, trzecia z pełnej wierności. Naturalność w $T$ daje żądany wzór.

Dla morfizmu $g:X\to Y$ w $\mathcal X$ sprzężenie i~\eqref{cat:eq:infty-topos-yoneda} identyfikują $\operatorname{Map}_{\mathcal X}(L(y(c)),g)$ z $(ig)(c)$. Jeżeli pierwsze morfizmy są równoważnościami dla wszystkich $c$, to $ig$ jest równoważnością przez twierdzenie~\ref{cat:thm:infty-presheaf-presentable}, a $g$ przez pełną wierność $i$. To dowodzi ostatniej tezy.
\end{proof}

\begin{przyklad}[Toposy presnopów]\label{cat:ex:infty-topos-presheaf}
Dla każdej małej $\mathcal C$ kategoria $\mathcal P(\mathcal C)$ jest $\infty$-toposem: za $L$ i $i$ bierzemy identyczność, która jest dostępna i zachowuje wszystkie granice. W szczególności $\mathcal S\simeq\mathcal P(*)$ jest $\infty$-toposem. Kategoria końcowa jest także $\infty$-toposem, gdyż $\mathcal P(\varnothing)$ jest kategorią funktorów z pustej kategorii i ma dokładnie jeden obiekt z ściągalną przestrzenią endomorfizmów. Pierwszy przykład pozwala badać zjawiska homotopijne już bez dodatkowej topologii.
\end{przyklad}

\begin{przyklad}[Dlaczego nie wystarcza dowolna dostępna lokalizacja]\label{cat:ex:infty-topos-truncation-localization}
W przykładzie~\ref{cat:ex:infty-bousfield-truncation-not-lex} rozpatrywaliśmy refleksję przestrzeni na przestrzenie dyskretne przez $\pi_0$. Przypomnijmy, że dla diagramu $*\to S^1\leftarrow *$ z oboma morfizmami wybierającymi ten sam punkt homotopijny iloczyn włóknisty ma typ przestrzeni pętli $\Omega S^1$, a $\pi_0(\Omega S^1)\cong\mathbb Z$. Po zastosowaniu $\pi_0$ do obiektów diagramu otrzymujemy diagram $*\to *\leftarrow *$, którego iloczyn włóknisty w zbiorach jest punktem. Porównanie nie jest bijekcją, więc ta refleksja nie zachowuje skończonych granic. Nie stanowi zatem prezentacji $\infty$-toposu w sensie definicji~\ref{cat:def:infty-topos}. Sam ten argument dotyczy konkretnego funktora lokalizacji; nie jest ogólnym kryterium wykluczającym inne prezentacje kategorii docelowej.
\end{przyklad}

\begin{uwaga}[Dalszy zakres]\label{cat:rem:infty-topos-scope}
Następnie wprowadzimy sita i snopy przestrzeni, warunek zstępowania dla nerwów Čecha oraz lewostronną dokładność snopifikacji. Potem omówimy aksjomaty Girauda, morfizmy geometryczne, obiekty obcięte i spójne, efektywne epimorfizmy, hiperprzykrycia i hiperzupełnienie. Rozróżnimy zwykły warunek snopu od warunku zstępowania względem wszystkich hiperprzykryć oraz powiążemy teorię z klasycznymi toposami i stosami.
\end{uwaga}
